% SEGMENT OLP-0584-S01
\documentclass[../../../include/open-logic-section]{subfiles}

\begin{document}

\olfileid{sth}{cardinals}{classing}
\olsection{முடிவுறு, \usetoken{S}{enumerable}, \usetoken{S}{nonenumerable}}

கண அளவெண்களை அறிமுகப்படுத்தியுள்ளதால், அவற்றின் வெவ்வேறு வகைகளைப்
பற்றிச் சிறிது பேசுவது பயனுள்ளது; குறிப்பாக முடிவுறு, !!{enumerable},
!!{nonenumerable} கண அளவெண்களைப் பற்றிப் பார்ப்போம்.

முதல் இரண்டு முடிவுகளிலிருந்து, முடிவுறு கண அளவெண்கள் துல்லியமாக
முடிவுறு வரிசையெண்களே என்று கிடைக்கும். அவற்றையே
\olref[z][infinity-again]{defnomega} இல் நமது \emph{இயல் எண்கள்}
என்று வரையறுத்தோம்:

\begin{prop}\ollabel{finitecardisoequal}
$n, m \in \omega$ என்க. அப்போது $n = m$ எனும்போதும் அப்போதுதான்
$\cardeq{n}{m}$.
\end{prop}

\begin{proof}
\emph{இடமிருந்து வலம்} செல்லும் திசை உடனடியாகத் தெரியும்.
\emph{வலமிருந்து இடம்} செல்லும் திசைக்கு, $n \neq m$ என்றாலும்
$\cardeq{n}{m}$ எனக் கொள்வோம். முக்கூற்றுப் பண்பின்படி $n \in m$
அல்லது $m \in n$; பொதுத்தன்மை இழப்பின்றி $n \in m$ எனக் கொள்வோம்.
அப்போது $n \subsetneq m$; மேலும் $f \colon m \to n$ என்ற
!!a{bijection} உள்ளது. எனவே $m$ டெடெகிண்ட் முடிவுறாதது.
இது \olref[z][infinity-again]{naturalnumbersarentinfinite} க்கு முரண்.
\end{proof}

\begin{cor}\ollabel{naturalsarecardinals}
$n \in \omega$ எனில், $n$ ஒரு கண அளவெண்.
\end{cor}

\begin{proof}
உடனடியாகப் பெறப்படுகிறது.
\end{proof}
\noindent

% SEGMENT OLP-0584-S02
ஒரு கண அளவெண்ணை ``முடிவுறு'' அல்லது ``முடிவுறாதது'' என்று
சொல்வதற்கான பல நியாயமான பொருள்கள் ஒன்றுபடுகின்றன என்றும் கிடைக்கிறது:
\begin{thm}\ollabel{generalinfinitycharacter}எந்தக் கணம் $A$ க்கும்,
பின்வருவன சமானமானவை:
\begin{enumerate}
	% SOURCE CORRECTION TA-STH-030: A அல்ல, அதன் கண அளவெண் இயல் எண் அல்ல.
	\item\ollabel{card:notinomega} $\card{A} \notin \omega$;
	அதாவது, $A$ இன் கண அளவெண் ஓர் இயல் எண் அல்ல;
	\item\ollabel{card:omegaplus} $\omega \leq \card{A}$;
	\item\ollabel{card:infinite} $A$ டெடெகிண்ட் முடிவுறாதது.
\end{enumerate}
\end{thm}

\begin{proof}
\olref[ord-arithmetic][using-addition]{ordinfinitycharacter},
\olref[cardinals][cardsasords]{lem:CardinalsBehaveRight},
\olref{naturalsarecardinals} ஆகியவற்றிலிருந்து பெறப்படுகிறது.
\end{proof}

இதனால் \olref[sfr][][]{part} இல் ஓரளவு முறைசாராமல் பயன்படுத்திய
சில கருத்துகளைப் பின்வருமாறு \emph{வரையறுக்க} முடிகிறது:

\begin{defn}\ollabel{defnfinite}
$\card{A}$ ஓர் இயல் எண், அதாவது $\card{A} \in \omega$,
எனும்போதும் அப்போதுதான் $A$ \emph{முடிவுறு} என்று சொல்கிறோம்.
இல்லாவிட்டால், $A$ \emph{முடிவுறாதது} என்று சொல்கிறோம்.
\end{defn}
\noindent
ஆனால், இந்த வரையறை $\ZFC$ இன் பின்னணியிலேயே அளிக்கப்படுகிறது
என்பதைக் கவனிக்கவும். எல்லாக் கணங்களுக்கும் கண அளவெண் இருப்பதை
உறுதிசெய்ய நன்கு வரிசைப்படுத்தல் தேவைப்பட்டது. உண்மையில், நன்கு
வரிசைப்படுத்தல் இல்லாவிட்டால், முடிவுறாததும் டெடெகிண்ட்
முடிவுறாததுமல்லாத ஒரு கணம் இருக்கக்கூடும். இவ்வகைச் சிக்கலுக்கு
\olref[choice][]{chap} இல் திரும்புவோம். இப்போது நன்கு
வரிசைப்படுத்தலைத் தொடர்ந்து பயன்படுத்துவோம்.

% SEGMENT OLP-0584-S03
இப்போது முடிவுறு கண அளவெண்களிலிருந்து முடிவுறாத கண அளவெண்களுக்குச்
செல்வோம். இரண்டு அடிப்படை முடிவுகள் இவை:

\begin{cor}\ollabel{omegaisacardinal}
$\omega$ தான் மீச்சிறிய முடிவுறாத கண அளவெண்.
\end{cor}

\begin{proof}
$\omega$ ஒரு கண அளவெண். ஏனெனில் $\omega$ டெடெகிண்ட்
முடிவுறாதது; ஏதேனும் $n \in \omega$ க்கு
$\cardeq{\omega}{n}$ எனில் $n$ உம் டெடெகிண்ட் முடிவுறாததாகும்.
இது \olref[z][infinity-again]{naturalnumbersarentinfinite} க்கு முரண்.
இப்போது, வரையறையின்படி $\omega$ தான் மீச்சிறிய முடிவுறாத கண அளவெண்.
\end{proof}

\begin{cor}
ஒவ்வொரு முடிவுறாத கண அளவெண்ணும் ஓர் எல்லை வரிசையெண்.
\end{cor}

\begin{proof}
$\alpha$ ஒரு முடிவுறாத தொடர்ச்சி வரிசையெண் என்க. அப்போது,
ஏதேனும் $\beta$ க்கு $\alpha = \beta \ordplus 1$.
% SOURCE CORRECTION TA-STH-031: முடிவுறு இயல் எண்ணின் தொடர்ச்சி மீண்டும் முடிவுறுவதே இங்கு தேவை.
$\beta$ முடிவுறாதது: அது ஓர் இயல் எண் எனில் அதன்
தொடர்ச்சியும் இயல் எண்களின் கணத்தில் இருக்கும்; இது முரண். முடிவுறு
வரிசையெண்களை வேறுபடுத்தும் \olref{finitecardisoequal} உடன் இதை
ஒப்பிடலாம். இப்போது
\olref[ord-arithmetic][using-addition]{ordinfinitycharacter} இன்படி
$\cardeq{\beta}{\beta \ordplus 1}$. எனவே
\olref[cardinals][cardsasords]{lem:CardinalsBehaveRight} இன்படி
$\card{\beta} = \card{\beta\ordplus 1} = \card{\alpha}$;
ஆகவே $\alpha \neq \card{\alpha}$.
\end{proof}

% SEGMENT OLP-0584-S04
\olref[sfr][siz][enm-alt]{defn:enumerable} இலேயே
!!{enumerable} முடிவுறாத கணங்களையும் !!{nonenumerable}
முடிவுறாத கணங்களையும் வேறுபடுத்தலாம் என்று குறிப்பிட்டோம்.
அந்த வரையறையிலிருந்து இயல்பாகப் பின்வருவது:

\begin{prop}
$A$ !!{enumerable} எனும்போதும் அப்போதுதான் $\card{A} \leq \omega$;
$A$ !!{nonenumerable} எனும்போதும் அப்போதுதான்
$\omega < \card{A}$.
\end{prop}

\begin{proof}
முக்கூற்றுப் பண்பின்படி இரு கூற்றுகளும் சமானமானவை. எனவே
$A$ !!{enumerable} எனும்போதும் அப்போதுதான்
$\card{A} \leq \omega$ என்பதை நிறுவினால் போதும்.
\emph{வலமிருந்து இடம்} செல்லும் திசையில், $\card{A} \leq \omega$
எனில், \olref[cardinals][cardsasords]{lem:CardinalsBehaveRight},
\olref{omegaisacardinal} ஆகியவற்றின்படி $\cardle{A}{\omega}$.
\emph{இடமிருந்து வலம்} செல்லும் திசையில், $A$ !!{enumerable} எனக்
கொள்வோம். அப்போது \olref[sfr][siz][enm-alt]{defn:enumerable} இன்படி
மூன்று நிலைகள் உண்டு:
\begin{enumerate}
	\item $A = \emptyset$ எனில்,
	\olref{naturalsarecardinals},
	\olref[cardinals][cardsasords]{lem:CardinalsBehaveRight}
	ஆகியவற்றின்படி $\card{A} = 0 \in \omega$.
	\item $\cardeq{n}{A}$ எனில்,
	\olref{naturalsarecardinals},
	\olref[cardinals][cardsasords]{lem:CardinalsBehaveRight}
	ஆகியவற்றின்படி $\card{A} = n \in \omega$.
	\item $\cardeq{\omega}{A}$ எனில்,
\olref{omegaisacardinal} இன்படி $\card{A} = \omega$.
\end{enumerate}
எனவே எல்லா நிலைகளிலும் $\card{A} \leq \omega$.
\end{proof}
\noindent
உண்மையில், $\omega$ க்கு ஒரு சிறப்பிடம் உள்ளது.
எண்ணக்கூடிய வரிசையெண்கள் பல இருந்தாலும்:

\begin{cor}
$\omega$ மட்டுமே !!{enumerable} முடிவுறாத கண அளவெண்.
\end{cor}

\begin{proof}
$\cardfont{a}$ ஓர் !!a{enumerable} முடிவுறாத கண அளவெண் என்க.
$\cardfont{a}$ முடிவுறாதது என்பதால் $\omega \leq \cardfont{a}$.
$\cardfont{a}$ ஓர் !!a{enumerable} கண அளவெண் என்பதால்,
$\cardfont{a} = \card{\cardfont{a}} \leq \omega$.
எனவே முக்கூற்றுப் பண்பின்படி $\cardfont{a} = \omega$.
\end{proof}

% SEGMENT OLP-0584-S05
நிச்சயமாக, முடிவுறாத எண்ணிக்கையில் கண அளவெண்கள் உள்ளன. எனவே
\emph{கண அளவெண்கள் எத்தனை?} என்று கேட்கலாம். அக்கேள்வியை
மறுபரிசீலனை செய்ய வேண்டும் என்பதைப் பின்வரும் முடிவுகள் காட்டுகின்றன.

\begin{prop}\ollabel{unioncardinalscardinal}
$X$ இன் ஒவ்வொரு உறுப்பும் ஒரு கண அளவெண் எனில், $\bigcup X$ உம்
ஒரு கண அளவெண்.
\end{prop}

\begin{proof}
$\bigcup X$ ஒரு வரிசையெண் என்பதை எளிதில் சரிபார்க்கலாம்.
$\alpha \in \bigcup X$ ஒரு வரிசையெண் என்க. அப்போது ஏதேனும்
கண அளவெண் $\cardfont{b}$ க்கு
$\alpha \in \cardfont{b} \in X$. $\cardfont{b}$ ஒரு கண அளவெண்
என்பதால் $\cardless{\alpha}{\cardfont{b}}$.
$\cardfont{b} \subseteq \bigcup X$ என்பதால்,
$\cardle{\cardfont{b}}{\bigcup X}$. எனவே
$\cardneq{\alpha}{\bigcup X}$. இதைப் பொதுமைப்படுத்தினால்,
$\bigcup X$ ஒரு கண அளவெண்.
\end{proof}

\begin{thm}\ollabel{lem:NoLargestCardinal}
மீப்பெரு கண அளவெண் என்று எதுவும் இல்லை.
\end{thm}

\begin{proof}
எந்தக் கண அளவெண் $\cardfont{a}$ க்கும், காண்டோரின் தேற்றமும்
(\olref[sfr][siz][car]{thm:cantor})
\olref[cardinals][cardsasords]{lem:CardinalsExist} உம் சேர்ந்து
$\cardfont{a} < \card{\Pow{\cardfont{a}}}$ என்பதைத் தருகின்றன.
\end{proof}

\begin{thm}
எல்லாக் கண அளவெண்களையும் கொண்ட கணம் இல்லை.
\end{thm}

\begin{proof}
மறுதலித்து நிறுவ, $C = \Setabs{\cardfont{a}}{\cardfont{a} \text{
ஒரு கண அளவெண்}}$ எனக் கொள்வோம். இப்போது
\olref{unioncardinalscardinal} இன்படி $\bigcup C$ ஒரு கண அளவெண்.
எனவே \olref{lem:NoLargestCardinal} இன்படி
$\cardfont{b} > \bigcup C$ ஆகுமாறு ஒரு கண அளவெண்
உள்ளது. வரையறையின்படி $\cardfont{b} \in C$; ஆகவே
$\cardfont{b} \subseteq \bigcup{C}$, எனவே
$\cardfont{b} \leq \bigcup C$. இது முரண்.
\end{proof}

இதை ரசலின் முரண்பாட்டுடனும் புராலி--ஃபோர்ட்டி முரண்பாட்டுடனும்
ஒப்பிட்டுப் பாருங்கள்.

\end{document}
